\documentclass[10pt]{article}
\usepackage[a4paper,margin=27mm]{geometry}
\usepackage{amsmath,amssymb,amsthm,mathtools,microtype}
\usepackage[colorlinks=true,linkcolor=blue,citecolor=blue,urlcolor=blue]{hyperref}
\hypersetup{pdfauthor={Volkov Sergey, SLAI},
 pdftitle={An explicit gap for three-point weighted selection in planar vector-valued regression}}
\newtheorem{theorem}{Theorem}
\newtheorem{lemma}[theorem]{Lemma}
\newcommand{\R}{\mathbb R}
\newcommand{\tr}{\operatorname{tr}}
\newcommand{\adj}{\operatorname{adj}}
\newcommand{\conv}{\operatorname{conv}}
\newcommand{\supp}{\operatorname{supp}}
\newcommand{\Fw}{F_{\mathrm{weighted}}(2,2,3)}
\newcommand{\T}{\mathcal T_3}
\title{An explicit gap for three-point weighted selection\\in planar vector-valued regression}
\author{Volkov Sergey\\SLAI}
\date{10 October 2026}
\begin{document}
\maketitle
\begin{abstract}
We prove $\Fw\le11/6-10^{-20}$ for minimum-Frobenius-norm
least-squares regression with two input and two output coordinates.
The argument combines a scalar certificate with quantitative stability of
centered planar selection and a seven-atom support reduction.
Together with the known lower bound, this gives
$13/8\le\Fw\le11/6-10^{-20}$.
Lean verifies the scalar selection theorem, its anisotropy improvement,
and exact numerical implications; the geometric argument is proved here.
\end{abstract}

\section{Problem and result}
Weighted data selection for regression was proposed by
Hanneke, Moran, Shlimovich and Yehudayoff~\cite{open,erms}.
Zhang~\cite{zhang}, Theorem~6.4, established
$13/8\le\Fw\le15/8$ and reduced the conjectured value $13/8$
to a finite planar moment problem. We improve the upper bound.

For a finite dataset $D=((x_i,y_i))_{i=1}^N\subset\R^2\times\R^2$, let
$L_D(W)=N^{-1}\sum_i\|Wx_i-y_i\|^2$ and $L_D^*=\min_W L_D(W)$.
For nonnegative weights $c$ of sum one, let $\widehat W_c$ be the
minimum-Frobenius-norm minimizer of $\sum_i c_i\|Wx_i-y_i\|^2$.
Then $\Fw$ is the supremum over datasets of
\[
 \frac{\inf_{|\supp c|\le3}L_D(\widehat W_c)}{L_D^*},
\]
with $0/0=1$. The infimum, including limits of full-rank selections,
is essential.

A residual frame is a finite collection $p_i>0$, $u_i,w_i\in\R^2$ with
\begin{equation}\label{frame}
 \sum_i p_i=1,\quad\|u_i\|^2=2,\quad
 \sum_i p_i u_i u_i^\top=I_2,\quad\sum_i p_i w_i u_i^\top=0.
\end{equation}
Put $E=\sum_i p_i\|w_i\|^2$ and $R=\sum_i p_iw_iw_i^\top$.
For independent $u_i,u_j$, let $G_{ij}u_i=w_i$, $G_{ij}u_j=w_j$.
Using Frobenius distance, define
\[
 \T=\min_{\substack{2\le|S|\le3\\S\text{ contains an independent pair}}}
 \operatorname{dist}\left(0,
 \conv\{G_{ij}:i,j\in S,\ \det(u_i,u_j)\ne0\}\right)^2.
\]
Tightness guarantees an independent pair.
\begin{theorem}\label{main}
Every residual frame satisfies $\T\le(5/6-10^{-20})E$.
Consequently $\Fw\le11/6-10^{-20}$.
\end{theorem}

\section{Planar selection and the scalar certificate}
\begin{lemma}\label{planar}
Let $v_j\in\R^2$ have positive probabilities $p_j$ and mean zero. For any
real $2\times2$ matrix $A$ with $\tr A=1$, some convex combination $v$ of
at most two of these points satisfies
\[
 \|v\|^2\le\sum_jp_j\|Av_j\|^2.
\]
\end{lemma}
\begin{proof}
Minimize the right-hand linear objective over centered probability vectors
on the fixed points, and choose a minimizer $q$ of least support. If the
columns $(1,v_j)$ on its positive support are dependent, a supported kernel
vector has both signs. Small perturbations in either direction remain
feasible, so optimality forces zero objective slope. Continuing until one
weight vanishes contradicts least support. Hence at most three affinely
independent points remain and their objective is no larger than the original.
If at most two remain, their centered mean itself is zero.

Otherwise let $S_q=\sum_jq_jv_jv_j^\top$. It is positive definite: a null
direction would place all three points on a line. Put $n_j=S_q^{-1}v_j$,
choose $i$ maximizing $r=\|n_i\|^2>0$, and set
\[
 \lambda_j=q_j\left(1-\frac{n_j^\top n_i}{r}\right).
\]
Cauchy--Schwarz and maximality give $\lambda_j\ge0$. Centering gives
$\sum_j\lambda_j=1$, and $\lambda_i=0$. Since
$\sum_jq_jv_jn_j^\top=I_2$,
\[
 v=\sum_j\lambda_jv_j=-\frac{n_i}{r},\qquad \|v\|^2=\frac1r.
\]
Finally, writing $C_q=\sum_jq_j\|Av_j\|^2$,
\[
 1=\sum_jq_j n_j^\top Av_j,\qquad
 1\le\left(\sum_jq_j\|n_j\|^2\right)C_q\le rC_q.
\]
Thus $\|v\|^2\le C_q$, proving the assertion.
\end{proof}

For positive semidefinite $S$ with $\tr S>0$, the matrix
$A=\adj S/\tr S$ has trace one, and in dimension two
\begin{equation}\label{adjidentity}
 \tr(ASA^\top)=\frac{\det S}{\tr S}.
\end{equation}
This identity holds also when $S$ is singular. If $\tr S=0$, all
positive-weight points vanish. Thus singular covariances cause no gap.

Fix an anchor $i$ and a unit vector $n_i\perp u_i$. Put
\[
 M_i=\tfrac12w_iu_i^\top,\quad
 b_{ij}=u_j^\top n_i,\quad c_{ij}=\tfrac12u_i^\top u_j,
 \quad \xi_{ij}=\frac{w_j-c_{ij}w_i}{b_{ij}}\quad(b_{ij}\ne0).
\]
Because $u_j=c_{ij}u_i+b_{ij}n_i$,
\[
 G_{ij}=M_i+\xi_{ij}n_i^\top,\qquad
 \|M_i+\xi n_i^\top\|_F^2=\tfrac12\|w_i\|^2+\|\xi\|^2.
\]
The cross term vanishes since $u_i^\top n_i=0$. The positive star weights
$P_{ij}=p_jb_{ij}^2$ satisfy
\[
 \sum_{b_{ij}\ne0}P_{ij}=1,\qquad
 \sum_{b_{ij}\ne0}P_{ij}\xi_{ij}=0.
\]
Indeed these are respectively $n_i^\top I_2n_i=1$ and the difference
between the cross-moment $\sum_jp_jw_ju_j^\top n_i=0$ and
$w_i u_i^\top I_2n_i/2=0$.

Define
\[
 L_i=\sum_{b_{ij}=0}p_j(w_j-c_{ij}w_i)(w_j-c_{ij}w_i)^\top\succeq0,
 \qquad H_i=R+\tfrac12w_iw_i^\top.
\]
The star covariance is
\begin{equation}\label{starcov}
 S_i=\sum_{b_{ij}\ne0}P_{ij}\xi_{ij}\xi_{ij}^\top=H_i-L_i.
\end{equation}
To check it, expand the same sum over all $j$. The two cross terms vanish
because $\sum_jp_jc_{ij}w_j=0$, and $\sum_jp_jc_{ij}^2=1/2$.
Subtracting the parallel terms yields \eqref{starcov}.

Apply Lemma~\ref{planar} with $A_i=\adj H_i/\tr H_i$; its denominator
$E+\|w_i\|^2/2$ is positive for $E>0$. PSD monotonicity and
\eqref{adjidentity} give a two-neighbor convex mean with
\[
 \|\xi\|^2\le\tr(A_iS_iA_i^\top)
 =\frac{\det H_i}{\tr H_i}-\tr(A_iL_iA_i^\top).
\]
Consequently, for every anchor,
\begin{equation}\label{anchor}
 \T\le \frac{\|w_i\|^2}{2}
 +\frac{\det(R+w_iw_i^\top/2)}{E+\|w_i\|^2/2}
 -\tr(A_iL_iA_i^\top).
\end{equation}
This involves at most three atoms. Dropping the last, nonnegative term is
valid but potentially loses information.

Normalize $E=1$ and rotate output coordinates so $R=\operatorname{diag}(r,1-r)$.
Set $a_i=w_{i1}^2$, $b_i=w_{i2}^2$, and $s=a+b$. The first two terms of
\eqref{anchor} are
\[
 g_r(a,b)=\frac{s}{2}+
 \frac{r(1-r)+((1-r)a+rb)/2}{1+s/2},\qquad
 \sum_i p_ia_i=r,\quad\sum_i p_ib_i=1-r.
\]
Put $\delta=2r-1$, $x=1+3\delta/10$, $y=1-3\delta/10$,
$P=xy$, $Q=7/3-21\delta^2/100$, and $h=P-ay-bx$.
Since
\[
 \sum_i p_i h(a_i,b_i)=\frac{21}{100}\delta^2\ge0,
\]
some positive-weight atom has $h\ge0$. Define
\[
 F(a,b)=s^2+\frac43s-\delta(a-b)-\frac73-\delta^2.
\]
Exact expansion gives
\begin{align}
 PF&=-\frac{121}{100}\delta^2P
     -h(Q+ax+by)-ab(x-y)^2,\label{poly}\\
 4(1+s/2)(g_r(a,b)-5/6)&=F(a,b).\label{gidentity}
\end{align}
Here $x,y,P,Q>0$ for $|\delta|\le1$. At the selected atom all terms on the
right of \eqref{poly} are nonpositive. This proves $\T\le5/6$.
Homogeneity gives $\T\le5E/6$ for general energy.

We also need its quantitative anisotropy consequence. At $h\ge0$,
$s\le P/\min(x,y)=\max(x,y)\le13/10$. Thus \eqref{poly}--\eqref{gidentity}
give
\begin{equation}\label{anisogap}
 \T\le\frac56-\frac{11}{60}\delta^2.
\end{equation}
The coefficient is $(121/100)/(4(1+13/20))=11/60$.


The scalar relaxation alone is sharp: the one-atom scalar distribution
$r=a=b=1/2$ has $g=5/6$. It is not a valid tight frame.
The strict improvement must therefore use geometry.

\section{Support reduction and quantitative stability}
\begin{lemma}\label{seven}
For each positive-energy frame there is a residual frame on a subset of
at most seven of its fixed atoms, with energy $E'\le E$, and with
$\T'\ge\T$. If $E'=0$, then $\T=0$.
\end{lemma}
\begin{proof}
Minimize $\sum_i q_i\|w_i\|^2$ over nonnegative weights on the fixed
atoms satisfying the frame constraints. The symmetric input covariance gives three linear equations and the
cross-covariance gives four. Normalization follows by taking the trace
of the input equation, since $\|u_i\|^2=2$. Thus there are at most
seven independent linear constraints.
The feasible probability polytope is compact and nonempty. At a minimizing
weight vector of least support, a dependence of its active constraint
columns would permit both signs of small feasible perturbation. The energy
slope must be zero by optimality. Continuing to the first zero weight would
reduce support while preserving energy, a contradiction. The positive
support consequently has at most seven atoms.
Deleting atoms can only increase the selection minimum, giving $\T'\ge\T$.
If $E'=0$, all remaining $w_i$ vanish; tightness still supplies an independent
pair and its interpolator is zero. Then $\T' =0$, so $\T=0$.
\end{proof}

\begin{lemma}\label{stability}
For a finite centered planar probability distribution, let $\tr A=1$,
$C=\sum_jp_j\|Av_j\|^2\le1/2$, and let $m\ge1/4$ be the minimum
squared norm of a convex combination of at most two original points.
Set $\zeta=C-m\ge0$. For any original point $v_*$ with $\|v_*\|\le2$,
\[
 C-\|Av_*\|^2\le50\zeta+10\sqrt\zeta\le60\sqrt\zeta.
\]
\end{lemma}
\begin{proof}
Choose a centered least-support minimizer $q$ of the linear $A$-cost. Since
$m>0$, its support is a positive triangle. Put $S_q=\sum_jq_jv_jv_j^\top$,
$n_j=S_q^{-1}v_j$, $r=\max_j\|n_j\|^2$, $M=\sum_jq_j\|n_j\|^2$,
$\alpha=1/M$, and $C_q=\sum_jq_j\|Av_j\|^2$. The planar construction and
weighted Cauchy--Schwarz yield
\[
 m\le1/r\le\alpha\le C_q\le C.
\]
Consequently $r\le4$, $1/4\le\alpha\le1/2$, and
\[
 r-M\le1/m-1/\alpha\le\zeta/(m\alpha)\le16\zeta.
\]
For $e_j=Av_j-\alpha n_j$, the trace identity gives
$\sum_jq_j\|e_j\|^2=C_q-\alpha\le\zeta$.
Write $c_j=\|Av_j\|^2$. Expanding and using Cauchy--Schwarz,
\begin{align*}
 \sum_jq_j|c_j-\alpha^2r|
 &\le\alpha^2(r-M)+2\alpha\sqrt{r\zeta}+\zeta
 \le5\zeta+2\sqrt\zeta,\\
 |\alpha^2r-C|&\le\alpha^2(r-M)+(C-\alpha)\le5\zeta.
\end{align*}
Thus $\sum_jq_j|c_j-C|\le10\zeta+2\sqrt\zeta$.
The affine coordinates of $v_*$ in the triangle are
$\beta_j=q_j(1+n_j^\top v_*)$: they sum to one and reproduce $v_*$
by the covariance identities. Since $\|n_j\|\le2$,
$|\beta_j|\le5q_j$. Inserting $v_*$ with small positive weight while
subtracting its affine coordinates from $q$ is feasible. LP optimality
implies $\|Av_*\|^2\ge\sum_j\beta_jc_j$. Therefore
\[
 C-\|Av_*\|^2\le\left|\sum_j\beta_j(c_j-C)\right|
 \le50\zeta+10\sqrt\zeta.
\]
Finally $0\le\zeta\le C\le1/2<1$, so $\zeta\le\sqrt\zeta$.
No bound on omitted points, or on triangle vertices themselves, was assumed.
\end{proof}


\section{The explicit gap}
\begin{proof}[Proof of the frame assertion in Theorem~\ref{main}]
If $E=0$, all $w_i$ vanish and $\T=0$.
For $E>0$, Lemma~\ref{seven} and homogeneity reduce the proof to at most
seven atoms with $E=1$. Put $\epsilon=10^{-20}$ and suppose
$\T\ge5/6-\epsilon$. By \eqref{anisogap},
$\delta^2\le60\epsilon/11$. Set $\eta=2\epsilon+|\delta|$; then
$\eta<3/10^{10}$, since $60/11<25/4$.

For $s=\|w_i\|^2$ the anchor bound is at most
$f(s)+|\delta|/2$, where
\[
 f(s)=s/2+(1+s)/(4+2s),\quad f(1)=5/6,\quad
 1/2\le f'(s)=1/2+1/(2(s+2)^2)\le5/8\quad(s\ge0).
\]
Indeed $r(1-r)\le1/4$ and
$((1-r)a+rb)/2\le(1+|\delta|)s/4$.
It follows that every $s_i\ge1-\eta$, and
\begin{equation}\label{energycontrol}
 \sum_i p_i(s_i-1+\eta)=\eta.
\end{equation}
Hence if $p_i\ge1/28$, then $1-\eta\le s_i\le1+27\eta$.

Retain these heavy atoms. The discarded mass is at most $1/4$.
Writing $u_i/\sqrt2=(\cos\theta_i,\sin\theta_i)$, tightness says that
$t_i=(\cos2\theta_i,\sin2\theta_i)$ has total mean zero.
The retained conditional mean has norm at most $1/3$.
There are therefore heavy $i,j$ with $t_i^\top t_j\le-1/2$, or equivalently
$c^2=(u_i^\top u_j/2)^2\le1/4$.
To see this, if zero lies in the retained hull, represent it using at most
three circle points. An antipodal pair, or a cyclic angular gap in
$[2\pi/3,\pi]$, gives the assertion. If zero is outside the hull,
its nearest point lies on a boundary edge (a vertex has norm one).
For unit endpoints the nearest point of their segment is its midpoint;
the squared distance is $(1+t_i^\top t_j)/2\le1/9$.

For the star at this heavy anchor, write
$H=R+w_iw_i^\top/2$, $A=\adj H/\tr H$, $\kappa=\det H/\tr H$,
$S=H-L$, $C=\tr(ASA^\top)$, and let $m$ be the minimum two-point
squared mean norm. Put $g_i=s_i/2+\kappa$.
Lemma~\ref{planar} and the anchored certificates give
$m\le C\le\kappa$ and $\T\le s_i/2+m$. Thus
\[
 0\le\zeta=C-m\le\kappa-m\le g_i-\T
 \le\frac{135}{8}\eta+|\delta|/2+\epsilon\le19\eta,
 \qquad m\ge1/3-\epsilon-27\eta/2\ge1/4.
\]
Put $v_k=w_k/\|w_k\|$ for $k=i,j$, and
$H_0=I_2/2+v_iv_i^\top/2$, $A_0=\adj H_0/(3/2)$.
All subsequent estimates hold for $\eta\le10^{-3}$, which is satisfied.
Output diagonalization gives
$\|R-I_2/2\|_{\mathrm{op}}=|\delta|/2$. Hence
\[
 \|H-H_0\|_{\mathrm{op}}\le14\eta,\quad
 |\tr H-3/2|\le27\eta/2,\quad \tr H\ge1.
\]
$H$ is positive definite, $A$ has trace one, and
$\kappa\le\tr H/4\le1/2$, so $C\le1/2$.
Writing $D=H-H_0$, the two-dimensional determinant expansion gives
\[
 |\det H-1/2|\le(3/2)\|D\|_{\mathrm{op}}+\|D\|_{\mathrm{op}}^2
 \le22\eta.
\]
It follows that $|\kappa-1/3|\le27\eta$. The identity
$\adj H/\tr H=I_2-H/\tr H$ gives
$\|A-A_0\|_{\mathrm{op}}\le23\eta$.

For the selected neighbor, put $b=b_{ij}$. Then
\[
 \xi=(w_j-cw_i)/b,\qquad \xi_0=(v_j-cv_i)/b,\qquad
 b^2=2(1-c^2)\ge3/2.
\]
The heavy norm bounds imply $\|\xi\|\le2$, $\|\xi_0\|\le3/2$,
and $\|\xi-\xi_0\|\le41\eta$; use
$|\sqrt{s_k}-1|\le|s_k-1|\le27\eta$ and $|c|\le1/2$.
As $A$ is positive definite of trace one, its operator norm is at most one.
Thus
\[
 \|A\xi-A_0\xi_0\|\le87\eta,\qquad
 |\|A\xi\|^2-\|A_0\xi_0\|^2|\le305\eta.
\]
For $d=v_i^\top v_j$, the eigenvalues of $H_0^{-1}$ are $1,2$, so
\[
 \|A_0\xi_0\|^2=
 \frac{(d-c)^2+4(1-d^2)}{18(1-c^2)}\le\frac{26}{81}.
\]
Indeed the numerator minus $(52/9)(1-c^2)$ equals
$(64/9)c^2-16/9-3(d+c/3)^2\le0$.
Consequently
\[
 \kappa-\|A\xi\|^2\ge1/81-332\eta\ge1/162.
\]
Here $\eta<3/10^{10}$.
Since $\kappa-C\le19\eta$, we obtain
\[
 C-\|A\xi\|^2\ge1/162-19\eta\ge1/200.
\]
All hypotheses of Lemma~\ref{stability} hold for this original star point.
It forces
\[
 1/200\le60\sqrt\zeta,\qquad \zeta\ge1/144000000.
\]
But $\zeta\le19\eta<57/10^{10}<1/144000000$, since
$57\cdot144000000=8208000000<10^{10}$. This contradiction proves
$\T<5/6-10^{-20}$ on the normalized reduced support, and the non-strict
bound transfers to all finite frames by Lemma~\ref{seven}.
\end{proof}


\section{Return to regression}
Let $W^\circ$ be the full-data minimum-norm minimizer and
$\rho_i=W^\circ x_i-y_i$. The normal equations imply
$\sum_i\rho_ix_i^\top=0$, and for every $B$,
\[
 L_D(W^\circ+B)=L_D^*+N^{-1}\sum_i\|Bx_i\|^2.
\]
Suppose $K=\sum_i x_ix_i^\top\succ0$.
Write $v_i=K^{-1/2}x_i$ and $\ell_i=\|v_i\|^2$.
For $\ell_i>0$, set
\[
 p_i=\ell_i/2,\qquad u_i=\sqrt{2/\ell_i}\,v_i,\qquad
 w_i=\sqrt{2/\ell_i}\,\rho_i.
\]
These satisfy \eqref{frame}, with $E=\sum_{\ell_i>0}\|\rho_i\|^2\le NL_D^*$.
For $\Theta=(W-W^\circ)K^{1/2}$, the excess loss is $\|\Theta\|_F^2/N$.

Given frame weights $q$, choose original weights proportional to $q_i/\ell_i$.
The selected objective, in $\Theta$ coordinates, is a positive multiple of
$\sum_iq_i\|\Theta u_i+w_i\|^2$. If $U(q)=\sum_iq_i u_iu_i^\top\succ0$,
its unique minimizer is $\Theta=-G(q)$, where
\begin{equation}\label{mixture}
 G(q)=\left(\sum_iq_iw_iu_i^\top\right)U(q)^{-1}
 =\sum_{i<j:D_{ij}>0}\frac{q_iq_jD_{ij}}{\sum_{a<b}q_aq_bD_{ab}}G_{ij},
 \qquad D_{ij}=\det(u_i,u_j)^2.
\end{equation}
This follows by expanding the determinant and adjugate in dimension two:
diagonal terms vanish and the $i,j$ numerator term is $D_{ij}G_{ij}$.

The hulls defining $\T$ are closures of full-rank fits on the same support
(cf.\ Zhang~\cite{zhang}, Lemma~E.1).
For three independent pairs, positive prescribed coefficients $\beta$
are attained by
$q_i:q_j:q_k=D_{jk}/\beta_{jk}:D_{ik}/\beta_{ik}:D_{ij}/\beta_{ij}$;
boundary coefficients follow by limits. For an anchored segment with
coefficients $\alpha_j$, take
\[
 q_i=1-t,\qquad q_j=t\frac{\alpha_j/D_{ij}}{\sum_k\alpha_k/D_{ik}},
 \qquad 0<t<1.
\]
Star-pair coefficients dominate the leaf pair as $t\downarrow0$, so
$G(q)\to\sum_j\alpha_jG_{ij}$. This also covers parallel leaves;
endpoints use an independent pair. Every approximating fit is full rank.
Consequently
\[
 \inf_{|\supp c|\le3}L_D(\widehat W_c)
 \le L_D^*+\T/N\le(11/6-10^{-20})L_D^*.
\]
No invariance of minimum-norm tie breaking under whitening is needed.

If the feature span has dimension one, write $x_i=a_ie$ and
$W^\circ=\beta e^\top$. On nonzero $a_i$, the planar vectors
$\rho_i/a_i$ are centered with probabilities $a_i^2/\sum_j a_j^2$.
Represent zero by at most three of them and choose original weights
proportional to the representation coefficients divided by $a_i^2$.
The selected normal equation recovers $\beta$; minimum norm sets the action
on $e^\perp$ to zero. Rank zero is immediate.
If $E=0$, a feature basis of size at most two recovers $W^\circ$ on the
feature span, including the zero-loss case. These cases have ratio one.
This proves Theorem~\ref{main}. The lower bound $13/8$ is
Zhang's Proposition~6.3~\cite{zhang}.

\section{Formal verification}
Source and Lean code are available at
\url{https://github.com/nZiben/weighted-selection-gap}.
The file \texttt{lean/Scalar.lean} verifies the two polynomial identities,
finite scalar moment selection, the bound $5/6-(11/60)(2r-1)^2$,
sharpness of the scalar relaxation, the ratio bound $26/9$, and the
exact numerical implications used above. The scalar selection theorems
assume only nonnegativity, normalization and the two stated moments;
they do not assume a selection bound.
The final numerical contradiction is explicitly conditional on the
geometric estimates proved in this paper. The planar selection and
stability lemmas, support reduction, and regression bridge are not
formalized; thus this is a partial formalization, not a Lean proof of
Theorem~\ref{main}.
The checked source uses Lean~4.34.0 and a pinned Mathlib revision.
The repository contains the compiler record and reproduction commands.
The reported theorem axioms are only \texttt{propext},
\texttt{Classical.choice} and \texttt{Quot.sound}; no proof placeholders
or additional axioms are used.

\section*{Use of AI}
The initial proof was produced by GPT-6 Pro (GPT-6 Astra) in the ChatGPT user interface.

\begin{thebibliography}{9}
\bibitem{open}
S. Hanneke, S. Moran, A. Shlimovich and A. Yehudayoff.
\emph{Open Problem: Data Selection for Regression Tasks}.
COLT 2025, PMLR 291, pp.\ 6225--6229.
\url{https://proceedings.mlr.press/v291/hanneke25e.html}.
\bibitem{erms}
S. Hanneke, S. Moran, A. Shlimovich and A. Yehudayoff.
\emph{Data Selection for ERMs}. 2025.
\url{https://arxiv.org/abs/2504.14572}.
\bibitem{zhang}
G. Zhang.
\emph{Exact Recovery Thresholds for Weighted Data Selection in Vector-Valued Linear Regression}.
2026, version 2 (2 October 2026).
\url{https://arxiv.org/abs/2608.30254v2}.
\end{thebibliography}
\end{document}
